\documentclass[11pt]{amsart}
\usepackage[a4paper,margin=1in]{geometry}
\usepackage[T1]{fontenc}
\usepackage{lmodern}
\usepackage{amsmath,amssymb,mathtools}
\usepackage{booktabs}
\usepackage{hyperref}
\usepackage{microtype}
\hypersetup{colorlinks=true,linkcolor=blue,citecolor=blue,urlcolor=blue}
\newtheorem{theorem}{Theorem}
\newtheorem{lemma}[theorem]{Lemma}
\newtheorem{remark}[theorem]{Remark}
\DeclareMathOperator{\sinc}{sinc}
\newcommand{\R}{\mathbb R}
\newcommand{\Z}{\mathbb Z}
\newcommand{\dd}{\,d}
\newcommand{\one}{\mathbf 1}
\newcommand{\Ctarget}{0.38056070}

\title{A Sharper Certified Lower Bound for Erd\H{o}s' Minimum-Overlap Problem}
\author{Deng Haowen}
\email{denghaowen59@gmail.com}
\date{7 September 2026}

\begin{document}
\begin{abstract}
We prove the computer-assisted lower bound
\[
  c_E>0.38056070
\]
for Erd\H{o}s' minimum-overlap constant.  The proof uses the 172-bin
mean-conditioned certificate framework of Price, retaining his certified Arb
bounds on the 170 noncentral bins and replacing the two binding center bins.
The new center certificate strengthens the earlier Parseval hybrid by replacing
a hard spectral cutoff with a fixed weighted Parseval row
\[
 -\sum_{n=1}^{400}w_n C_{n\pi}\le\frac12,
 \qquad 0\le w_n\le1,
\]
and by adding exact harmonic rows $C_{n\pi}\le0$.  The frozen certificate has
one second-moment multiplier, 80 ordinary cosine multipliers, one weighted
Parseval multiplier, and 216 harmonic multipliers.  Two independent interval
implementations verify the same frozen proof object.  A direct MPFR calculation
gives the rigorous upper bound $2.62770156549654007831\ldots$ for
$\int q_+$; an independent \texttt{mpmath.iv} verifier gives
\[
 \int_{-2}^{2}q_+(t)\dd t
 \le 2.62770156515276293334556455703
 <\frac1{0.38056070}.
\]
The optimizer used to discover the certificate is outside the trusted proof
path.
\end{abstract}
\maketitle

\section{Continuous formulation}
For a partition $A\sqcup B=\{1,\ldots,2N\}$ with $|A|=|B|=N$, put
\[
 r_{A,B}(m)=|\{(a,b)\in A\times B:a-b=m\}|,
 \qquad R(A,B)=\max_m r_{A,B}(m),
\]
and
\[
 c_E=\liminf_{N\to\infty}\frac1N
 \min_{A\sqcup B=\{1,\ldots,2N\}}R(A,B).
\]
Following White and Price \cite{White2023,Price2026}, let $\mathcal H$ be the
measurable functions $h:[0,2]\to[0,1]$ with $\int_0^2h=1$, put $g=1-h$ on
$[0,2]$, extend both by zero, and define
\[
 p_h(t)=\int_{\R}h(x)g(x+t)\dd x,\qquad -2\le t\le2.
\]
Then $p_h\ge0$, $\int p_h=1$, and $\operatorname{supp}p_h\subset[-2,2]$.
Write
\[
 c_{\rm cont}=\inf_{h\in\mathcal H}\|p_h\|_\infty.
\]

\begin{lemma}[Discrete-to-continuous reduction]\label{lem:dc}
$c_E\ge c_{\rm cont}$.
\end{lemma}
\begin{proof}
For $I_k=[(k-1)/N,k/N)$ set
$h_N=\sum_{a\in A}\one_{I_a}$ and $g_N=\sum_{b\in B}\one_{I_b}$.
With $\tau(u)=(1-|u|)_+$,
\[
 p_{h_N}(t)=\frac1N\sum_{a\in A,b\in B}\tau(Nt-(b-a)).
\]
For fixed $x$, the nonzero values among $\tau(x-d)$, $d\in\Z$, sum to at most
one.  Grouping by $d=b-a$ gives
$\|p_{h_N}\|_\infty\le R(A,B)/N$.  Minimize and take the liminf.
\end{proof}

Fix $h\in\mathcal H$, write $p=p_h$, and let
\[
 \mu=\int_{-2}^{2}t p(t)\dd t.
\]
The elementary rearrangement bound $1/2\le\mathbb E X\le3/2$ for a random
variable with density $h$ implies $\mu\in[-1,1]$.

\section{Analytic rows}
\subsection{Second moment}
\begin{lemma}\label{lem:m2}
For every admissible $p$,
\[
 \int t^2p(t)\dd t=\frac23+\frac{\mu^2}{2}.
\]
Hence on $|\mu|\le1/320$,
\[
 \int t^2p(t)\dd t\le B_2:=\frac23+\frac1{2\cdot320^2}.
\]
\end{lemma}
\begin{proof}
Let $X,Y$ be independent with densities $h,g$.  Then $p$ is the density of
$Y-X$, while
$\mathbb EX+\mathbb EY=2$ and
$\mathbb EX^2+\mathbb EY^2=8/3$.  Since
$\mathbb EY-\mathbb EX=\mu$, one has
$\mathbb EX=1-\mu/2$, $\mathbb EY=1+\mu/2$, and the displayed identity follows.
\end{proof}

\subsection{Fourier rows}
Define
\[
 H(\xi)=\int_0^2h(x)e^{i\xi(x-1)}\dd x,
 \qquad G(\xi)=\int_0^2g(x)e^{i\xi(x-1)}\dd x,
\]
and
\[
 C_\xi=\int_{-2}^{2}p(t)\cos(\xi t)\dd t.
\]
Since $H(\xi)+G(\xi)=2\sinc\xi$ and
$\int p(t)e^{i\xi t}\dd t=\overline{H(\xi)}G(\xi)$, the usual completion of
squares gives the following.

\begin{lemma}[Pointwise cosine bound]\label{lem:cos}
For every real $\xi$,
\[
 C_\xi\le\sinc^2\xi.
\]
In particular, for every positive integer $n$,
\[
 C_{n\pi}=-|H(n\pi)|^2\le0.
\]
\end{lemma}

\subsection{Weighted Parseval}
The next form of White's Parseval restriction is the only new ingredient in the
certificate language relative to the previous hard-cutoff version.

\begin{lemma}[Weighted Parseval row]\label{lem:wp}
Let $(w_n)$ be finitely supported with $0\le w_n\le1$.  Then
\[
 -\sum_{n\ge1}w_n C_{n\pi}\le\frac12.
\]
\end{lemma}
\begin{proof}
At $n\pi$, Lemma~\ref{lem:cos} gives
$-C_{n\pi}=|H(n\pi)|^2$.  Parseval for $u\mapsto h(u+1)$ on $[-1,1]$ gives
\[
 1+2\sum_{n\ge1}|H(n\pi)|^2\le2,
\]
because $H(0)=1$ and $0\le h\le1$ implies
$\int h^2\le\int h=1$.  Therefore
\[
 -\sum_nw_nC_{n\pi}
 =\sum_nw_n|H(n\pi)|^2
 \le\sum_n|H(n\pi)|^2\le\frac12.
\]
\end{proof}
\begin{remark}
The underlying Parseval energy inequality is due to the existing Fourier
framework and is not claimed as new.  The present improvement is the use of a
frozen non-rectangular spectral window inside Price's mean-conditioned dual
certificate.
\end{remark}

\section{Dual certificate}
\begin{lemma}[Positive-part dual bound]\label{lem:dual}
Suppose $p$ satisfies finitely many valid inequalities
$\int p\phi_j\le B_j$.  For $\lambda_j\ge0$ put
\[
 q(t)=1+\sum_j\lambda_j(B_j-\phi_j(t)),
 \qquad q_+(t)=\max(q(t),0).
\]
Then
\[
 \|p\|_\infty\ge\frac1{\int_{-2}^{2}q_+(t)\dd t}.
\]
\end{lemma}
\begin{proof}
$\int pq\ge1$, while $p\ge0$ implies
$\int pq\le\|p\|_\infty\int q_+$.
\end{proof}

\section{The frozen center certificate}
Price's 172-bin cover has the binding pair
\[
 I_-=[-1/320,0],\qquad I_+=[0,1/320].
\]
The machine-readable proof object
\path{certificate/weighted_center_certificate_038056070.json}
replaces precisely these two bins.  It contains exact decimal strings for
\begin{itemize}
\item one nonnegative second-moment multiplier $\lambda_2$;
\item 80 nonnegative ordinary cosine multipliers $\lambda_\xi$;
\item one nonnegative weighted-Parseval multiplier $\lambda_W$ and 400 fixed
weights $0\le w_n\le1$;
\item 216 nonnegative exact-harmonic multipliers $\eta_n$.
\end{itemize}
The associated even trigonometric polynomial is
\begin{align*}
q(t)=1&+\lambda_2(B_2-t^2)
+\sum_{\xi\in\Xi}\lambda_\xi(\sinc^2\xi-\cos(\xi t))\\
&+\lambda_W\left(\frac12+\sum_{n=1}^{400}w_n\cos(n\pi t)\right)
-\sum_{n\in\mathcal N}\eta_n\cos(n\pi t).
\end{align*}
All rows are valid on both center bins by Lemmas
\ref{lem:m2}, \ref{lem:cos}, and \ref{lem:wp}.  The optimizer used to find the
coefficients is not needed to check any of these facts.

\section{Two independent interval verifications}
The release contains two implementations that read or embed the same frozen
proof object but use different interval libraries and subdivision strategies.
Floating root searches, when used, only propose subdivisions and have no
logical role.

\subsection{Direct MPFR path}
\texttt{code/verify\_weighted\_center\_mpfr.c} uses MPFR with explicit directed
rounding and no interval library shared with the Python implementation.  At a
cell center it encloses $q,q',\ldots,q^{(6)}$ and uses a global seventh-derivative
bound for the Taylor remainder.  Cells proved positive are integrated with the
explicit elementary antiderivative of $q$; terminal ambiguous cells are charged
conservatively by their width times a rigorous upper bound for $q$.  The 32
archived decimal subintervals cover $[0,2]$ exactly, and the runner parses each
printed upward-rounded bound as an exact rational before summing and doubling
by evenness.  A full reproduction gives
\begin{equation}\label{eq:mpfr}
 \int_{-2}^{2}q_+(t)\dd t
 \le 2.627701565496540078311925225883\ldots
 <\frac1{0.38056070}.
\end{equation}
The direct-MPFR path reads a compact terminating-decimal coefficient transcription
derived from the frozen JSON proof object.
\path{code/check_weighted_mpfr_certificate_match.py} verifies that transcription
against the JSON certificate before the release check is accepted.

\subsection{Independent \texttt{mpmath.iv} path}
\texttt{code/verify\_weighted\_center\_mpmath.py} independently reconstructs
all ordinary Fourier right-hand sides using \texttt{mpmath.iv}.  Its 32-piece
manifest uses third-order Taylor bounds on the main mass and sixth-order Taylor
bounds in the high-frequency tail.  Exact rational aggregation gives
\begin{equation}\label{eq:mpmath}
 \int_{-2}^{2}q_+(t)\dd t
 \le 2.62770156515276293334556455703.
\end{equation}
Since
\[
 \frac1{0.38056070}
 =2.62770170435360246079009209306\ldots,
\]
the certified $D$-side margin in this second implementation is
$1.392008395\times10^{-7}$.  Thus every admissible $p$ with
$\mu\in I_-\cup I_+$ satisfies
\[
 \|p\|_\infty>0.38056070.
\]

\section{Noncentral bins and the theorem}
For the other 170 mean bins we reuse Price's published Arb-certified upper
bounds \cite{Price2026}.  The largest such value is attained at bins 77 and 94:
\[
 D_{\rm noncentral}
 \le2.6275385308733757900902721758784\ldots,
\]
so
\[
 \frac1{D_{\rm noncentral}}
 \ge0.3805843333028524\ldots>0.38056070.
\]
The release script \texttt{code/check\_noncentral\_target.py} parses all 170
vendored Arb balls and checks this comparison exactly from their printed
midpoint-radius enclosures.

\begin{theorem}\label{thm:main}
The Erd\H{o}s minimum-overlap constant satisfies
\[
 \boxed{c_E>0.38056070}.
\]
\end{theorem}
\begin{proof}
Every mean $\mu\in[-1,1]$ belongs to Price's 172-bin cover.  On the two center
bins, \eqref{eq:mpmath} and Lemma~\ref{lem:dual} give the stated bound.  On every
other bin the reused Arb certificate gives the stronger noncentral bound above.
Hence $c_{\rm cont}>0.38056070$, and Lemma~\ref{lem:dc} implies the theorem.
\end{proof}

\section{Reproducibility and trust model}
The trusted finite data are the frozen JSON certificate, the vendored Price Arb
reports for the 170 noncentral bins, and the source code of the two interval
verifiers.  The LP search and floating-point continuous integration are
exploratory only.

The release provides:
\begin{itemize}
\item exact certificate validation, including $0\le w_n\le1$ and nonnegative
multipliers;
\item a direct-MPFR verifier and exact chunk aggregator;
\item an independent \texttt{mpmath.iv} verifier with a separate manifest and
exact rational aggregation;
\item a target check for all 170 noncentral Arb reports;
\item SHA-256 hashes for the release files and a reproducible PDF build.
\end{itemize}
The two implementations were developed within the same project; this is not a
third-party audit.  External reproduction remains desirable.

\paragraph{Use of AI.}
Computational exploration, certificate search, verifier engineering, and parts
of the exposition were carried out with substantial assistance from OpenAI's
ChatGPT.  The language model is not part of the trusted proof path.  The human
author is responsible for the mathematical claims and release.

\begin{thebibliography}{9}
\bibitem{White2023}
E. P. White, \emph{A new bound for Erd\H{o}s' minimum overlap problem},
Acta Arith. \textbf{208} (2023), 235--255.
\href{https://doi.org/10.4064/aa220728-7-6}{doi:10.4064/aa220728-7-6}.
\bibitem{Price2026}
L. Price, \emph{A Certified Lower Bound for Minimum Overlap}, proof note and
reproducible Arb certificate for $c_E>0.38055470$, 29 June 2026.
\url{https://github.com/Leeham06972452/erdos-36-lower-bound}.
\bibitem{Occisn2026}
occisn, \emph{Independent verification of the $0.38055470$ minimum-overlap certificate}, 2026.
\url{https://github.com/occisn/erdos-36-certified-lower-bound}.
\end{thebibliography}
\end{document}
